\documentclass[10pt,a4paper]{amsart}
\usepackage[margin=19mm]{geometry}
\usepackage{amsmath,amssymb,mathtools}
\usepackage[scr=rsfs,cal=euler]{mathalfa}
\usepackage{xfrac}
\usepackage[unicode,hidelinks,bookmarks=false]{hyperref}
\newcommand{\ie}{i.e.\ }
\newcommand{\cf}{cf.\ }
\newcommand{\resp}{respectively\ }
\newcommand{\Subsection}{\subsection}
\providecommand{\nobreakdash}{\nobreak\hbox{-}}
\setlength{\emergencystretch}{2em}
\multlinegap0pt
\usepackage{bm}
\usepackage{bbm}
\usepackage{dsfont}
\usepackage{xspace}
\usepackage{cancel}
\usepackage{mathtools}
\usepackage[shortlabels]{enumitem}
\usepackage{amsthm}
\makeatletter
\@ifundefined{lemma}{\newtheorem{lemma}{Lemma}}{}
\makeatother
\newcommand{\DeltaS}{\Delta_{\src}}
\newcommand{\DeltaT}{\Delta_{\trg}}
\newcommand{\Ebb}{\mathbb E}
\newcommand{\KL}{\mathrm{KL}}
\newcommand{\R}{\mathbb R}
\newcommand{\Xsrc}{X^{(\src)}}
\newcommand{\Xtar}{X^{(\trg)}}
\newcommand{\ip}[2]{\langle #1,#2\rangle}
\newcommand{\nS}{n_{\src}}
\newcommand{\nT}{n_{\trg}}
\newcommand{\sigmaS}{\sigma_{\src}}
\newcommand{\sigmaT}{\sigma_{\trg}}
\newcommand{\sphere}[1]{\mathbb S^{#1}}
\newcommand{\src}{S}
\newcommand{\trg}{T}                   
\newcommand{\zsrc}{z^{(\src)}}
\newcommand{\ztar}{z^{(\trg)}}
\title[Transfer Learning in High-Dimensional Clustering: Minimax Th]{Transfer Learning in High-Dimensional Clustering: Minimax Thresholds and Applications in Single-Cell Data}
\author{Abhinav Chakraborty, Sagnik Nandy}
\date{Source: arXiv:2607.25031v1 (CC BY 4.0)}
\begin{document}
\maketitle

\noindent\emph{Source and licence.} Transfer Learning in High-Dimensional Clustering: Minimax Thresholds and Applications in Single-Cell Data. Version arXiv:2607.25031v1 (CC BY 4.0), distributed under the Creative Commons Attribution 4.0 International licence, as recorded in the arXiv licence metadata for this exact version and shown on the abstract page https://arxiv.org/abs/2607.25031v1. Full rights notice: licenses/arxiv-2607.25031-CC-BY-4.0.txt.\par\smallskip

\noindent\textit{Editorial scope.} Counted unit: the Lemma whose source label is \texttt{lem:sec3\_product\_row\_kl} in the article (source0.tex lines 4466--4478, source label \texttt{lem:sec3\_product\_row\_kl}), delivered together with the article's own complete original proof of it (source0.tex lines 4480--4628, one \texttt{proof} environment of 149 lines). No mathematics is written, repaired or re-derived: statement and proof are the article's own text line for line. Required context retained uncounted: eq:def-Xi\_k (lines 4155--4157); eq:def-l-defs (lines 4162--4166); eq:def-rT-rS (lines 4168--4172); eq:def-xr-yr (lines 4174--4178); lem:sec3\_product\_row\_density (lines 4203--4217); eq:sec3\_product\_row\_density (lines 4208--4216); lem:sec3\_product\_odd\_pert (lines 4245--4265); lem:sec3\_product\_moment\_bds (lines 4316--4345); lem:sec3\_product\_approx\_error (lines 4368--4382); lem:sec3\_product\_score\_sq (lines 4425--4436). Every definition, display and label of those lines is retained unchanged. Omitted from this package and not redistributed: 1--4154, 4158--4161, 4167--4167, 4173--4173, 4179--4202, 4217--4244, 4266--4315, 4346--4367, 4383--4424, 4437--4465, 4479--4479, 4629--7065. Nothing inside a retained excerpt is omitted or altered: every retained body is delivered exactly as the article's lines are written, so no omission receipt of any kind accompanies this package. Citations: the retained excerpts carry no citation command at all, so no bibliography entry of the article is transcribed and no reference is redistributed. Reference adaptation: none was needed and none was made; every reference inside the delivered unit resolves inside it, so no unresolved reference or citation is produced. Printed slips kept literal: no printed word, symbol, hypothesis or number of the counted statement or of its proof was altered. Layout and class: the article's own class is \texttt{article} with packages including geometry, xr, inputenc, amsmath, amssymb, amsthm, mathabx, mathtools, bm, booktabs, bbm, enumitem, natbib, hyperref; the wrapper is the lane's pinned \texttt{amsart} editorial wrapper at 10pt on A4, which restates the theorem-like environments the retained text needs and adds the article's own macro definitions that the retained text needs (\texttt{\textbackslash DeltaS}, \texttt{\textbackslash DeltaT}, \texttt{\textbackslash Ebb}, \texttt{\textbackslash KL}, \texttt{\textbackslash R}, \texttt{\textbackslash Xsrc}, \texttt{\textbackslash Xtar}, \texttt{\textbackslash ip}, \texttt{\textbackslash nS}, \texttt{\textbackslash nT}, \texttt{\textbackslash sigmaS}, \texttt{\textbackslash sigmaT}, \texttt{\textbackslash sphere}, \texttt{\textbackslash src}, \texttt{\textbackslash trg}, \texttt{\textbackslash zsrc}, \texttt{\textbackslash ztar}), with extra wrapper packages (bm, bbm, dsfont, xspace, cancel, mathtools, enumitem). The delivered numbering is the unit's own. Assets: no retained span contains a figure environment, a graphics-inclusion command, a PDF-page inclusion, a table or a TikZ picture; no image file is copied into this package, the package declares no asset dependency and no image file is redistributed with it. Licence: version arXiv:2607.25031v1 of this article is distributed under the Creative Commons Attribution 4.0 International licence, as recorded in the arXiv licence metadata for this exact version and shown on the abstract page https://arxiv.org/abs/2607.25031v1. A full rights notice is delivered at licenses/arxiv-2607.25031-CC-BY-4.0.txt. Characters: every retained excerpt is pure ASCII, so the delivered engine, XeLaTeX with its default Latin Modern text font, reports no missing character.
\begin{equation}\label{eq:def-Xi_k}
    \Xi_{j_0} := (\eta,\tau_{-j_0},\zsrc),
\end{equation}

\begin{equation}\label{eq:def-l-defs}
    l^{(\trg)}_\pm := \frac{\ztar_\pm}{\sqrt{\nT}}\in\sphere{\nT-1},
    \qquad
    l^{(\src)} := \frac{\zsrc}{\sqrt{\nS}}\in\sphere{\nS-1},
\end{equation}

\begin{equation}\label{eq:def-rT-rS}
    r_T := \DeltaT\sqrt{\frac{\nT}{d}},
    \qquad
    r_S := \DeltaS\sqrt{\frac{\nS}{d}}.
\end{equation}

\begin{equation}\label{eq:def-xr-yr}
    x_r := \left(\frac{\Xtar_{i,r}}{\sigmaT}\right)_{i=1}^{\nT}\in\R^{\nT},
    \qquad
    y_r := \left(\frac{\Xsrc_{j,r}}{\sigmaS}\right)_{j=1}^{\nS}\in\R^{\nS},
\end{equation}

\begin{lemma}[Conditional row density]
\label{lem:sec3_product_row_density}
With $\Xi_{j_0}$ as in~\eqref{eq:def-Xi_k}, $l^{(\trg)}_\pm$, $l^{(\src)}$ as in~\eqref{eq:def-l-defs}, $r_T$, $r_S$ as in~\eqref{eq:def-rT-rS}, and $x_r$, $y_r$, $a_\pm$, $b$ as in~\eqref{eq:def-xr-yr}, fix $\Xi_{j_0}$.
The conditional row law of $(x_r,y_r)$ under $\tau_{j_0}=\pm1$
has density relative to the standard Gaussian reference measure given by
\begin{equation}
\label{eq:sec3_product_row_density}
    p_{\pm,r}(x_r,y_r)
    =
    e^{-(r_T^2+r_S^2)/2}\Bigl[
        \cosh\!\bigl(a_\pm(x_r)\bigr)\cosh\!\bigl(b(y_r)\bigr)
        + \tilde\mu\sinh\!\bigl(a_\pm(x_r)\bigr)\sinh\!\bigl(b(y_r)\bigr)
    \Bigr].
\end{equation}
\end{lemma}

\begin{lemma}[Odd perturbation bounds for the row log-density]
\label{lem:sec3_product_odd_pert}
Define
\[
    h(a,b):=\log\!\bigl(\cosh(a)\cosh(b)+\tilde\mu\sinh(a)\sinh(b)\bigr).
\]
Set $\Delta_1:=\partial_a h(a,b)-\tanh(a)$ and
$\Delta_2:=\partial_{aa}h(a,b)-\operatorname{sech}^2(a)$.  Then for all
$a,b\in\R$:
\begin{enumerate}[label=(\roman*)]
\item $\displaystyle\partial_a h(a,b)
      =
      \frac{\tanh(a)+\tilde\mu\tanh(b)}{1+\tilde\mu\tanh(a)\tanh(b)}$,
      and $|\partial_a h(a,b)|\le 1$.
\item $\partial_{aa}h(a,b)=1-(\partial_a h(a,b))^2$,
      and $0\le\partial_{aa}h(a,b)\le 1$.
\item $|\partial_b\Delta_1|\le\tilde\mu$ and $|\Delta_1|\le\tilde\mu|b|$.
\item $\Delta_2=-(2\tanh(a)+\Delta_1)\Delta_1$,
      so $|\Delta_2|\le 2\tilde\mu|b|+\tilde\mu^2 b^2$.
\end{enumerate}
\end{lemma}

\begin{lemma}[Moment bounds for the rotated coordinates]
\label{lem:sec3_product_moment_bds}
With $r_T$, $r_S$ as in~\eqref{eq:def-rT-rS} and $l^{(\src)}$ as in~\eqref{eq:def-l-defs}, fix
$l^{(\trg)}_+=e_1\in\sphere{\nT-1}$ and let $(x,y)\sim p_{+,r}$, where
$p_{+,r}$ is the row density~\eqref{eq:sec3_product_row_density} with
$a(x)=r_T X_1$ ($X_1$ the first coordinate of $x$) and
$b(y)=r_S\ip{l^{(\src)}}{y}$.
Let $\alpha\in\R$ and define
\[
    Y_\alpha:=\cos\alpha\,X_1+\sin\alpha\,X_2, \quad
    Z_\alpha:=-\sin\alpha\,X_1+\cos\alpha\,X_2, \quad
    B:=r_S\ip{l^{(\src)}}{y}.
\]
Assume $r_T\le1$.  Then there is a representation
\begin{align}
    X_1 &= r_T\eta+g_1, & \eta&\sim\mathrm{Rad}(1/2),\ g_1\sim N(0,1),
    \label{eq:sec3_product_X1_decomp}\\
    B   &= r_S^2 s\eta+r_S h, & s&\sim\mathrm{Rad}\!\bigl(\tfrac{1+\tilde\mu}{2}\bigr),\
    h\sim N(0,1), \label{eq:sec3_product_B_decomp}
\end{align}
where $(\eta,s,g_1,h)$ are independent, and $X_2=g_2\sim N(0,1)$ is
independent of $(X_1,B)$.  Uniformly in $\alpha$:
\begin{gather*}
    \Ebb|Y_\alpha|^j\le C\;\text{ for }j=1,2,3,4,6, \qquad
    \Ebb Z_\alpha^4\le C,\\
    \Ebb|B|^4\le C(r_S^8+r_S^4), \qquad
    \Ebb|r_T Y_\alpha|^4\le Cr_T^4, \qquad
    |\Ebb[Y_\alpha B]|\le C\tilde\mu r_T r_S^2.
\end{gather*}
\end{lemma}

\begin{lemma}[Linear approximation error for the row score]
\label{lem:sec3_product_approx_error}
With $r_T$, $r_S$ as in~\eqref{eq:def-rT-rS} and the notation of
Lemma~\ref{lem:sec3_product_moment_bds}, set $a:=r_T Y_\alpha$, $b:=B$,
and $q(a,b):=\partial_a h(a,b)$.  Assume $r_T\le c$ and
$\tilde\mu^2 r_S^2\le c$.  Then
\begin{equation}
\label{eq:sec3_product_approx_error}
    \bigl|\Ebb\bigl[Y_\alpha\bigl(q(a,b)-a-\tilde\mu b\bigr)\bigr]\bigr|
    \le
    C\bigl(r_T^3+\tilde\mu^2 r_T r_S^2\bigr),
\end{equation}
and consequently
$r_T\bigl|\Ebb[Y_\alpha(q-a-\tilde\mu b)]\bigr|\le C(r_T^4+\tilde\mu^2 r_T^2 r_S^2)$.
\end{lemma}

\begin{lemma}[Score-square moment bound]
\label{lem:sec3_product_score_sq}
With $r_T$, $r_S$ as in~\eqref{eq:def-rT-rS} and the notation of
Lemma~\ref{lem:sec3_product_moment_bds}, set $a:=r_T Y_\alpha$, $b:=B$,
and $q(a,b):=\partial_a h(a,b)$.  Assume $r_T\le c$.  Then
\begin{equation}
\label{eq:sec3_product_score_sq}
    r_T^2\,\Ebb\!\bigl[Z_\alpha^2\,q(r_T Y_\alpha,B)^2\bigr]
    \le
    C\bigl(r_T^4+\tilde\mu^2 r_T^2 r_S^2\bigr).
\end{equation}
\end{lemma}

\begin{lemma}[Conditional row KL bound]
\label{lem:sec3_product_row_kl}
With $\Xi_{j_0}$ as in~\eqref{eq:def-Xi_k}, $r_T$, $r_S$ as in~\eqref{eq:def-rT-rS}, and $l^{(\trg)}_\pm$ as in~\eqref{eq:def-l-defs},
for each fixed $\Xi_{j_0}$ let $p_{\pm,r}$ be the conditional row densities from
Lemma~\ref{lem:sec3_product_row_density}.  Whenever $r_T\le c$ and
$\tilde\mu^2 r_S^2\le c$,
\begin{equation}
\label{eq:sec3_product_row_kl}
    \KL(p_{+,r}\,\|\,p_{-,r})
    \le
    C\,r_T^2(r_T^2+\tilde\mu^2 r_S^2)\,\|l^{(\trg)}_+-l^{(\trg)}_-\|^2.
\end{equation}
\end{lemma}

\begin{proof}
\textit{Reduction by rotation.}
KL divergence is invariant under bijective measurable maps of the sample
space: for any bijection $T$, $\KL(P\|Q)=\KL(T_\#P\|T_\#Q)$.  We exploit
this with the map $T:(x_r,y_r)\mapsto(Ox_r,y_r)$, where $O\in O(\nT)$ is the
orthogonal transformation satisfying
\[
    Ol^{(\trg)}_+ = e_1, \qquad Ol^{(\trg)}_- = l_\alpha := \cos\alpha\,e_1+\sin\alpha\,e_2,
    \qquad \alpha := \arccos\!\ip{l^{(\trg)}_+}{l^{(\trg)}_-}\in[0,\pi].
\]
(Such an $O$ exists because $l^{(\trg)}_\pm$ are unit vectors; $l_\alpha$ is the
unique unit vector in $\mathrm{span}\{e_1,e_2\}$ with $\ip{e_1}{l_\alpha}=\ip{l^{(\trg)}_+}{l^{(\trg)}_-}$.)

\textit{Why only $x_r$ is rotated.}
The two hypotheses $\tau_{j_0}=\pm1$ differ only through the target statistic
$a_\pm(x_r)=r_T\ip{l^{(\trg)}_\pm}{x_r}$.
The source statistic $b(y_r)=r_S\ip{l^{(\src)}}{y_r}$ is the same under
both hypotheses (since $l^{(\src)}=\zsrc/\sqrt{\nS}$ is fixed by $\Xi_{j_0}$),
so no rotation of $y_r$ is needed or helpful.  We therefore act on
$x_r$ alone, leaving $y_r$ unchanged.

\textit{Effect of $T$ on the row distributions.}
The pushforward $T_\#P_{\pm,r}$ is the law of $(Ox_r,y_r)$ when
$(x_r,y_r)\sim P_{\pm,r}$.  Since $p_{\pm,r}$
(Lemma~\ref{lem:sec3_product_row_density}) depends on $x_r$ only through
$a_\pm(x_r)=r_T\ip{l^{(\trg)}_\pm}{x_r}$, and
$\ip{l^{(\trg)}_\pm}{x_r}=\ip{Ol^{(\trg)}_\pm}{Ox_r}$, setting
$\tilde x_r:=Ox_r$ shows that $T_\#P_{+,r}$ has target statistic
$r_T\ip{e_1}{\tilde x_r}$ and $T_\#P_{-,r}$ has target statistic
$r_T\ip{l_\alpha}{\tilde x_r}$, while both retain the same source statistic
$b(y_r)$.  Hence $T_\#P_{+,r}=P_{e_1}$ and $T_\#P_{-,r}=P_{l_\alpha}$,
where the subscripts now denote the (rotated) target label directions.
For the densities to be expressed against the same reference measure
$Q=N(0,I_{\nT})\otimes N(0,I_{\nS})$ after the change of variables, we
need $T_\#Q=Q$.  The source factor $N(0,I_{\nS})$ is unchanged since $T$
does not touch $y_r$.  The target factor $N(0,I_{\nT})$ is $O$-invariant
because $\|Ox_r\|=\|x_r\|$.  Therefore $T_\#Q=Q$, and
\[
    \KL(P_{+,r}\|P_{-,r})
    =
    \KL(T_\#P_{+,r}\|T_\#P_{-,r})
    =
    \KL(P_{e_1}\|P_{l_\alpha})
    =:
    F(\alpha).
\]

\textit{Definitions in rotated coordinates.}
Write $X_i:=\ip{e_i}{x_r}$ for the $i$-th coordinate of the (now rotated)
standardized column, and define (as in Lemma~\ref{lem:sec3_product_moment_bds})
\[
    Y_\alpha := \ip{l_\alpha}{x_r} = \cos\alpha\,X_1+\sin\alpha\,X_2,
    \qquad
    Z_\alpha := \partial_\alpha Y_\alpha = -\sin\alpha\,X_1+\cos\alpha\,X_2,
\]
so that $\partial_\alpha Z_\alpha=-Y_\alpha$.  After the rotation,
$a_+(x_r)=r_T X_1$ and $a_-(x_r)=r_T Y_\alpha$.  Also set
\[
    B := b(y_r) = r_S\ip{l^{(\src)}}{y_r},
\]
the source statistic from~\eqref{eq:def-xr-yr}, which is the same under
both hypotheses.  Because
$h(a,b)=\log(\cosh a\cosh b+\tilde\mu\sinh a\sinh b)$ is the log-density
of each row (Lemma~\ref{lem:sec3_product_row_density}), and the Gaussian
prefactor $e^{-(r_T^2+r_S^2)/2}$ cancels in the KL ratio,
\[
    F(\alpha)
    =
    \Ebb_{p_{e_1}}\!\bigl[h(r_T X_1,B)-h(r_T Y_\alpha,B)\bigr].
\]

\textit{Law under $p_{e_1}$.}
From the model (see proof of Lemma~\ref{lem:sec3_product_row_density}):
conditional on $\Xi_{j_0}$ and $\tau_{j_0}=+1$, the target column (after rotation) is distributed as
$x_r \overset{d}{=} \xi_r\,r_T\,e_1+\varepsilon$, where $\xi_r\sim\mathrm{Rad}(1/2)$
is the coordinate direction sign and $\varepsilon\sim N(0,I_{\nT})$ is independent.
Thus $X_2=\ip{e_2}{x_r}=\varepsilon_2\sim N(0,1)$ is independent of $(\xi_r,X_1,B)$.
More generally, $(Y_\alpha,Z_\alpha)$ is an orthogonal rotation of $(X_1,X_2)$, and
since under $p_{e_1}$ we have $\Ebb[X_1^2]=r_T^2+1$ and $\Ebb[X_2^2]=1$:
\begin{equation}\label{eq:rotated-moments}
    \Ebb[Y_\alpha^2]=r_T^2\cos^2\alpha+1, \qquad \Ebb[Z_\alpha^2]=r_T^2\sin^2\alpha+1.
\end{equation}

\textit{$F(0)=F'(0)=0$.}
$F(0)=0$ since $Y_0=X_1$.  Differentiating in $\alpha$ and writing
$q:=\partial_a h(r_T Y_\alpha,B)$ (Lemma~\ref{lem:sec3_product_odd_pert}(i)):
\[
    F'(\alpha) = -r_T\,\Ebb_{p_{e_1}}\!\bigl[q(r_T Y_\alpha,B)\,Z_\alpha\bigr].
\]
At $\alpha=0$, $Z_0=X_2$ is independent of $(X_1,B)$ and $\Ebb[X_2]=0$,
so $F'(0)=-r_T\,\Ebb[q(r_T X_1,B)]\cdot\Ebb[X_2]=0$.

\textit{Computing $F''$.}
Differentiating once more, using $\partial_\alpha Z_\alpha=-Y_\alpha$ and
$\partial_a q=\partial_{aa}h=1-q^2$ (Lemma~\ref{lem:sec3_product_odd_pert}(ii)):
\begin{align*}
    F''(\alpha)
    &=
    -r_T^2\Ebb\!\bigl[(1-q^2)Z_\alpha^2\bigr]
    +
    r_T\Ebb\!\bigl[q\,Y_\alpha\bigr]\\
    &=
    r_T\Ebb[q\,Y_\alpha]
    -r_T^2\Ebb[Z_\alpha^2]
    +r_T^2\Ebb[Z_\alpha^2 q^2].
\end{align*}
Writing $q = (r_T Y_\alpha+\tilde\mu B)+(q-r_T Y_\alpha-\tilde\mu B)$
in the first term $r_T\Ebb[qY_\alpha]$ and expanding gives the decomposition
\[
    F''(\alpha) = T_0+T_1+T_2+T_3,
\]
where
\begin{align*}
    T_0 &:= r_T^2\Ebb[Y_\alpha^2-Z_\alpha^2], &
    T_1 &:= \tilde\mu r_T\Ebb[Y_\alpha B], \\
    T_2 &:= r_T\Ebb\!\bigl[Y_\alpha(q-r_T Y_\alpha-\tilde\mu B)\bigr], &
    T_3 &:= r_T^2\Ebb[Z_\alpha^2 q^2].
\end{align*}

\textit{Bound on $T_0$.}
By~\eqref{eq:rotated-moments},
$\Ebb[Y_\alpha^2]-\Ebb[Z_\alpha^2]=r_T^2(\cos^2\alpha-\sin^2\alpha)$,
so $|T_0|=r_T^4|\cos^2\alpha-\sin^2\alpha|\le r_T^4$.

\textit{Bound on $T_1$.}
By Lemma~\ref{lem:sec3_product_moment_bds}, $|\Ebb[Y_\alpha B]|\le\tilde\mu r_T r_S^2$,
so $|T_1|=\tilde\mu r_T|\Ebb[Y_\alpha B]|\le\tilde\mu^2 r_T^2 r_S^2$.

\textit{Bound on $T_2$.}
This is the approximation error; Lemma~\ref{lem:sec3_product_approx_error} gives
$|T_2|\le C(r_T^4+\tilde\mu^2 r_T^2 r_S^2)$.

\textit{Bound on $T_3$.}
Lemma~\ref{lem:sec3_product_score_sq} gives
$|T_3|=r_T^2\Ebb[Z_\alpha^2 q^2]\le C(r_T^4+\tilde\mu^2 r_T^2 r_S^2)$.

\textit{Conclusion.}
Collecting the four bounds,
$|F''(\alpha)|\le Cr_T^2(r_T^2+\tilde\mu^2 r_S^2)$ uniformly in $\alpha$.
Since $F(0)=F'(0)=0$, Taylor's theorem in integral form gives
\[
    F(\alpha) = \int_0^\alpha(\alpha-t)F''(t)\,dt
    \le Cr_T^2(r_T^2+\tilde\mu^2 r_S^2)\alpha^2.
\]
Finally, $\|l^{(\trg)}_+-l^{(\trg)}_-\|^2=2(1-\cos\alpha)=4\sin^2(\alpha/2)$,
and $|\sin\theta|\ge 2|\theta|/\pi$ for $\theta\in[-\pi/2,\pi/2]$, so
$\alpha^2\le\frac{\pi^2}{4}\|l^{(\trg)}_+-l^{(\trg)}_-\|^2$,
proving~\eqref{eq:sec3_product_row_kl}.
\end{proof}

\end{document}
